\documentclass[12pt]{article}
\usepackage[spanish]{babel}
\usepackage[utf8]{inputenc}
\usepackage[T1]{fontenc}
\usepackage{amsmath, amssymb}
\usepackage{geometry}
\usepackage{enumitem}
\usepackage{needspace}
\usepackage{tikz}
\usetikzlibrary{babel}

\geometry{letterpaper, margin=2cm}
\setlength{\parindent}{0pt}
\setlength{\parskip}{6pt}

\newcommand{\puntografico}{0.09}
\newcommand{\puntocerrado}[2]{\fill (#1,#2) circle (\puntografico);}

\begin{document}

\begin{center}
  {\Large \textbf{Borrador de items}}\\
  {\large Bloque 2. Relaciones y Algebra - Afirmacion 2}\\[4pt]
  Prueba Nacional Estandarizada de Matematica, Secundaria 2026
\end{center}

\section*{Afirmacion}
Resuelve problemas relacionados con la inversa de una funcion en sus distintas representaciones.

\section*{Items propuestos}

\begin{enumerate}[label=\textbf{\arabic*.}, itemsep=1.05cm]

\Needspace{0.95\textheight}
\item
La siguiente representacion grafica corresponde a la funcion $p$, que determina la presion $p(t)$, en kilopascales, registrada por un sensor industrial en funcion del tiempo $t$, en minutos, desde que se inicia una prueba, con $0 \leq t \leq 8$.

\begin{center}
\begin{tikzpicture}[scale=0.72]
  \path[use as bounding box] (-1.1,-0.9) rectangle (9.2,7.1);
  \draw[<->, thick] (-0.55,0) -- (8.8,0);
  \draw[<->, thick] (0,-0.45) -- (0,6.65);
  \node[anchor=west] at (8.9,-0.05) {$t$};
  \node[anchor=south east] at (-0.05,6.7) {$p(t)$};
  \node[anchor=north] at (8,0) {$8$};
  \node[anchor=east] at (0,2) {$2$};
  \node[anchor=east] at (0,6) {$6$};
  \draw[densely dashed, gray!75] (8,0) -- (8,2) -- (0,2);
  \draw[densely dashed, gray!75] (0,6) -- (0,6);
  \draw[thick] (0,6) -- (8,2);
  \puntocerrado{0}{6}
  \puntocerrado{8}{2}
  \node[anchor=south] at (3.2,4.55) {$p$};
\end{tikzpicture}
\end{center}

De acuerdo con la informacion anterior, \textquestiondown cual representacion grafica corresponde a $p^{-1}$?

\begin{enumerate}[label=\Alph*)]
  \item
  \begin{tikzpicture}[scale=0.45, baseline=-0.5ex]
    \draw[<->, thick] (-0.4,0) -- (8.7,0);
    \draw[<->, thick] (0,-0.4) -- (0,6.7);
    \node[anchor=west] at (8.8,0) {$x$};
    \node[anchor=south] at (0,6.8) {$y$};
    \node[anchor=north] at (8,0) {$8$};
    \node[anchor=east] at (0,2) {$2$};
    \node[anchor=east] at (0,6) {$6$};
    \draw[densely dashed, gray!75] (8,0) -- (8,2) -- (0,2);
    \draw[thick] (0,6) -- (8,2);
    \fill (0,6) circle (0.12);
    \fill (8,2) circle (0.12);
  \end{tikzpicture}
  \item
  \begin{tikzpicture}[scale=0.45, baseline=-0.5ex]
    \draw[<->, thick] (-0.4,0) -- (8.7,0);
    \draw[<->, thick] (0,-0.4) -- (0,8.7);
    \node[anchor=west] at (8.8,0) {$x$};
    \node[anchor=south] at (0,8.8) {$y$};
    \node[anchor=north] at (2,0) {$2$};
    \node[anchor=north] at (6,0) {$6$};
    \node[anchor=east] at (0,8) {$8$};
    \draw[densely dashed, gray!75] (2,0) -- (2,0);
    \draw[densely dashed, gray!75] (6,0) -- (6,8) -- (0,8);
    \draw[thick] (2,0) -- (6,8);
    \fill (2,0) circle (0.12);
    \fill (6,8) circle (0.12);
  \end{tikzpicture}
  \item
  \begin{tikzpicture}[scale=0.45, baseline=-0.5ex]
    \draw[<->, thick] (-0.4,0) -- (8.7,0);
    \draw[<->, thick] (0,-0.4) -- (0,8.7);
    \node[anchor=west] at (8.8,0) {$x$};
    \node[anchor=south] at (0,8.8) {$y$};
    \node[anchor=north] at (2,0) {$2$};
    \node[anchor=north] at (8,0) {$8$};
    \node[anchor=east] at (0,6) {$6$};
    \draw[densely dashed, gray!75] (8,0) -- (8,6) -- (0,6);
    \draw[thick] (2,0) -- (8,6);
    \fill (2,0) circle (0.12);
    \fill (8,6) circle (0.12);
  \end{tikzpicture}
  \item
  \begin{tikzpicture}[scale=0.45, baseline=-0.5ex]
    \draw[<->, thick] (-0.4,0) -- (8.7,0);
    \draw[<->, thick] (0,-0.4) -- (0,8.7);
    \node[anchor=west] at (8.8,0) {$x$};
    \node[anchor=south] at (0,8.8) {$y$};
    \node[anchor=north] at (2,0) {$2$};
    \node[anchor=north] at (6,0) {$6$};
    \node[anchor=east] at (0,8) {$8$};
    \draw[densely dashed, gray!75] (2,0) -- (2,8) -- (0,8);
    \draw[densely dashed, gray!75] (6,0) -- (6,0);
    \draw[thick] (2,8) -- (6,0);
    \fill (2,8) circle (0.12);
    \fill (6,0) circle (0.12);
  \end{tikzpicture}
\end{enumerate}

\Needspace{0.56\textheight}
\item
Considere la siguiente informacion:

En un vivero se usa la funcion $a$, que determina la altura $a(s)$, en centimetros, de una planta de albahaca en funcion del tiempo $s$, en semanas, desde que se trasplanto, dada por
\[
a(s)=7s+18,
\]
donde $s$ representa el tiempo en semanas, con $2 \leq s \leq 9$.

De acuerdo con la informacion anterior, si se requiere determinar el criterio de la funcion inversa de $a$, entonces, \textquestiondown cual opcion corresponde a ese criterio?

\begin{enumerate}[label=\Alph*)]
  \item $s(a)=7a+18$
  \item $s(a)=\dfrac{a-18}{7}$
  \item $s(a)=\dfrac{a}{7}-18$
  \item $s(a)=\dfrac{18-a}{7}$
\end{enumerate}

\Needspace{0.62\textheight}
\item
Considere la siguiente informacion:

Una empresa de empaque establece que la funcion $c$, que determina el costo $c(h)$, en miles de colones, de operar una maquina segun el tiempo $h$, en horas, esta dada por
\[
c(h)=4(h^2+6),
\]
donde $h$ representa el tiempo de operacion de la maquina, con $3 \leq h \leq 7$.

De acuerdo con la informacion anterior, si se requiere determinar el criterio de la funcion $h$, la cual corresponde a la funcion inversa de $c$, entonces, \textquestiondown cual opcion corresponde a ese criterio?

\begin{enumerate}[label=\Alph*)]
  \item $h(c)=\sqrt{\dfrac{c}{4}-6}$
  \item $h(c)=\sqrt{\dfrac{c}{4}+6}$
  \item $h(c)=\sqrt{\dfrac{c-6}{4}}$
  \item $h(c)=\dfrac{\sqrt{c}}{4}-6$
\end{enumerate}

\Needspace{0.56\textheight}
\item
Considere la siguiente informacion:

La funcion $q$, que determina la cantidad de litros de solucion nutritiva $q(r)$ que hay en un deposito, esta dada por
\[
q(r)=5r+12,
\]
donde $r$ representa el tiempo en horas desde que inicia el llenado del deposito, con $4 \leq r \leq 14$.

De acuerdo con la informacion anterior, si se define la funcion $r$, la cual corresponde a la funcion inversa de $q$, entonces, \textquestiondown cual es el ambito de $r$?

\begin{enumerate}[label=\Alph*)]
  \item $[4,14]$
  \item $[17,82]$
  \item $[32,82]$
  \item $[32,70]$
\end{enumerate}

\Needspace{0.58\textheight}
\item
Una plataforma de cursos registra una relacion entre cada codigo de curso y la cantidad de horas de tutoria asignadas por semana. Para que esa relacion tenga funcion inversa, \textquestiondown cual condicion debe cumplirse?

\begin{enumerate}[label=\Alph*)]
  \item Cada codigo de curso debe estar asociado con una unica cantidad de horas, aunque dos codigos diferentes tengan la misma cantidad.
  \item Cada cantidad de horas debe estar asociada con al menos dos codigos de curso.
  \item Cada codigo de curso debe estar asociado con una unica cantidad de horas y no debe haber dos codigos diferentes con la misma cantidad.
  \item Cada codigo de curso debe estar asociado con todas las cantidades de horas disponibles.
\end{enumerate}

\end{enumerate}

\section*{Clave y justificacion breve}

\begin{enumerate}[label=\textbf{\arabic*.}]
  \item \textbf{D.} La grafica de la inversa se obtiene al intercambiar las coordenadas de los extremos de $p$: $(0,6)$ pasa a $(6,0)$ y $(8,2)$ pasa a $(2,8)$.
  \item \textbf{B.} Al despejar $a=7s+18$, se obtiene $s=\dfrac{a-18}{7}$.
  \item \textbf{A.} Al despejar $c=4(h^2+6)$, se obtiene $\dfrac{c}{4}=h^2+6$ y, como $h$ es positivo en el dominio dado, $h=\sqrt{\dfrac{c}{4}-6}$.
  \item \textbf{A.} El ambito de la funcion inversa corresponde al dominio de la funcion original. Como $4 \leq r \leq 14$, el ambito de la inversa es $[4,14]$.
  \item \textbf{C.} Para que una funcion tenga inversa como funcion, debe ser uno a uno: cada entrada tiene una unica imagen y valores de entrada diferentes no comparten la misma imagen.
\end{enumerate}

\end{document}
